% En este archivo se incluuye la info de la segunda paguina del TT, sepa como se llame XD

% Evitar la numeración en esta página
\thispagestyle{empty}


% Dibujamos encima de la página con coordenadas absolutas
\begin{tikzpicture}[remember picture,overlay]

    % Imagen en la esquina superior izquierda 
    \node at ([xshift=3.3cm,yshift=24.7cm]current page.south west) 
    {\includegraphics[width=5.5cm]{Logos/ipn.png}};
        
    % Imagen en la esquina superior derecha
    \node at ([xshift=18cm,yshift=24.7cm]current page.south west) 
    {\includegraphics[width=2.5cm]{Logos/escom.png}};

\end{tikzpicture}

 % Aquí va la información xd
    {\fontsize{16pt}{14pt}\selectfont
    \begin{center}
    \textbf{INSTITUTO POLITÉCNICO NACIONAL}\\
    \textbf{ESCUELA SUPERIOR DE CÓMPUTO}\\
    \textbf{SUBDIRECCIÓN ACADÉMICA}\\
    \end{center}
    }

    \vspace{0.5cm} % Ps pa dar espacio pa q más
    
    {\fontsize{12pt}{14pt}\selectfont 
    No de TT: 2026-B045  \hfill Fecha de presentación de TT-II: Mayo, 2026
    }
    
    \vspace{0cm} % Ps pa dar espacio pa q más

    {\fontsize{14pt}{14pt}\selectfont
    \begin{center}
    Documento técnico
    \end{center}
    }
    
    \vspace{0cm} % Ps pa dar espacio pa q más
    
    {\fontsize{18pt}{14pt}\selectfont
    \begin{center}
        \textbf{"Navegación con asistencia auditiva en tiempo real para personas con visión reducida, mediante el uso de lentes adaptados e inteligencia artificial, para el manejo de siniestros"}
    \end{center}
    }
    
    
    {\fontsize{14pt}{14pt}\selectfont
    \begin{center}
    \textit{Presentan}\\
    \bf {Thania Paola Esquivel García\footnote{tesquivelg2000@alumno.ipn.mx}}\\
    \bf{José Carlos Espinosa Martínez \footnote{jespinosam1900@alumno.ipn.mx}}
    \end{center}
    }
    
    {\fontsize{14pt}{14pt}\selectfont
    \begin{center}
    \textit{Director}\\
    \bf {M. en C. Ávila Sánchez Cristhian Alejandro }
    \end{center}
    }
    
  
    {\fontsize{12pt}{14pt}\selectfont
    \begin{center}
    \textbf{RESUMEN}\\
    \end{center}
    La vista es uno de los sentidos que nos ayudan a interpretar el entorno e interactuar de manera efectiva con él. Sin embargo, personas que carecen de manera total o parcial de este sentido se enfrentan a desafíos diarios que limitan su independencia y acceso a diversas áreas cotidianas. La navegación para personas cuya visión es reducida o presentan discapacidad visual es esencial para su autonomía y bienestar. Si bien existen métodos tradicionales de asistencia, estos suelen tener un enfoque  que no se ajusta a las necesidades de viajar a diversos puntos. Este trabajo terminal tiene como objetivo utilizar la inteligencia artificial para proporcionar asistencia auditiva en tiempo real a personas con visión reducida, mediante el uso de lentes adaptados, así como un sistema de navegación. El sistema permite la identificación del entorno para guiar al usuario de manera adecuada, promoviendo una mayor independencia y seguridad para los usuarios.\\
    \textbf{Palabras clave:} Discapacidad visual, Inteligencia artificial, Reconocimiento de voz, Visión artificial, Navegación asistida.}